\section{Architecture}\label{sec:architecture}

\begin{figure}[htbp]
\centering
\fig{architecture}
\caption{The three roles and the protocols between them.}
\end{figure}

Ymir has three roles, each an independent toggle in a node's config. One \code{ymird} process
can run all three, or they can be spread over any number of devices in any combination. Roles
always talk through the same protocols; when two share a device they use an in-process fast
path, so behaviour, auth and failure handling do not change with placement.

\begin{xltabular}{\linewidth}{@{}L{2.4cm}YY@{}}
\toprule
Role & Does & Never does \\
\midrule
\endhead
Model node & Runs engines, advertises hardware and models in a signed capability record, obeys Load and Unload, serves raw inference. & Routing, placement, conversation state, talking to apps. \\
Inference node & Builds the pool catalog, picks a model node per request, resumes failed answers, owns context sync, aggregates health and alerts. & Running models (unless the same device is also a model node). \\
App node & Finds inference nodes, reads the catalog, sends prompts. & Discovering or calling model nodes directly. \\
\bottomrule
\end{xltabular}

\begin{table}[htbp]
\centering\small
\begin{tabularx}{\linewidth}{@{}lY@{}}
\toprule
Between & Protocol \\
\midrule
App to inference & \code{/weirdnet-agent/infer/2} and \code{/weirdnet-agent/catalog/1} \\
Inference to model & \code{/weirdnet-agent/infer/2} (raw mode) and \code{/weirdnet-backends/node/1} \\
Inference to inference & \code{/weirdnet/sync/2} for context and policy; \code{/ymir/coord/1} for coordination \\
Any to any (admin) & \code{/ymir/control/1} for \code{ymirctl} commands \\
\bottomrule
\end{tabularx}
\end{table}

% ---------------------------------------------------------------------------
\subsection{Weirdnet transport and discovery}\label{sec:transport}

\begin{figure}[htbp]
\centering
\fig{transport-flow}
\caption{Finding a peer and being allowed to talk to it are separate steps.}
\end{figure}

In short: finding a peer and being allowed to talk to it are separate steps. The relay and mDNS
only help nodes locate each other; what admits a caller is the roster check that every serving
node runs on every call.

Every node and client is a Weirdnet node with a certificate chained to the owner's identity, and
Noise authenticates every stream end to end, relays included. Nodes find each other by key, by
the user record and by the apps each peer announces, with LAN hints and mDNS for same-network
peers. Presence never grants anything; the roster does. By default nodes bootstrap from the
public Weirdnet relay, and the operator can override that.

The \code{infer/2} protocol is CBOR throughout and streams tokens as frames with backpressure.
It adds priority classes (interactive, batch, background), a stable job id with an attempt
number, a conversation reference, an optional assistant prefix for resume, Queued frames from a
real priority queue, offset-carrying Delta and Reasoning frames, whole-only ToolCall frames,
embeddings, and response metadata (which node served it, cache hit, stand-in flag, resume
offsets). Dropping a call resets its stream and aborts the engine request, freeing the slot.
Large payloads travel as content-addressed roots rather than inline frames.

Discovery and remote inference are deliberately generic: they live in the published
\code{weirdnet-agent} crate, so any Weirdnet app can be an app node, and Bragi works with any
inference node that serves the catalog and infer protocols.

% ---------------------------------------------------------------------------
\subsection{Scheduling and placement}\label{sec:scheduling}

\begin{figure}[htbp]
\centering
\fig{request-path}
\caption{How one call is distributed.}
\end{figure}

Inference nodes route each request from the capability records model nodes push to them every
15 seconds and on change. Among nodes with the model ready, the built balancer scores each by
predicted completion and takes the cheaper of two random candidates (power of two choices), so
stale records do not make every coordinator pile onto one node.
\begin{align*}
  \text{cost} &= \text{wait} + \text{service} + \text{load} + \text{penalty}\\
  \text{wait} &= \frac{\text{queue\_depth}}{\text{slots}} \times \text{mean\_service}
\end{align*}

\begin{figure}[htbp]
\centering
\fig{cost-example}
\caption{\textsc{Illustrative} example, not a measurement. Node A: 2 queued, 2 slots, answered
Busy 3\,s ago. Node B: 2 queued, 1 slot, slower service.}
\end{figure}

Illustrative example, not a measurement: the two candidates the balancer drew, scored in seconds
with the formula in the code (\code{ymir-core::balance}). Node A has the shorter queue per slot
but is still carrying the Busy penalty, so the request goes to node B. Mean service starts at 2
seconds until a node has finished a request for that model; a cold node would add its load-time
estimate.

\code{queue_depth} is the larger of the node's reported in-flight count and what this
coordinator has itself sent and not seen finish; \code{mean_service} is measured per node and
model; a Busy answer adds a flat 4-second penalty to that node for the next 5 seconds; and a
node on the same device wins when two costs are within 5 percent. Model nodes keep the last word
through admission, answering Busy when slots and queue are full, so several coordinators can run
active-active without overcommitting a node.

Fairness is between devices and priority classes, since the pool has one owner. Interactive
requests are served first, and each device has a cap on batch and background work so one
laptop's batch job cannot starve a phone's chat. \code{ymirctl batch} fans a file of prompts out
across the pool at batch priority with per-task retries and results in order.

\begin{table}[htbp]
\centering\small
\begin{tabularx}{\linewidth}{@{}llY@{}}
\toprule
Class & Who uses it & How it is treated \\
\midrule
interactive & Bragi chat, the bridge & Served first; never held back by the device cap \\
batch & \code{ymirctl batch} & Yields between requests; counted against the device cap \\
background & Prefix warming, re-placement (designed) & Lowest; counted against the device cap \\
\bottomrule
\end{tabularx}
\end{table}

The design also specifies cloud-wide swap and placement: a scheduler lease held by one inference
node, a signed and synced placement plan, replication of popular models when queue wait exceeds
a target, prefix warming, and a pool-wide semantic cache. These are designed but not built yet (milestone M6) \planned.

% ---------------------------------------------------------------------------
\subsection{Failover and continuation}\label{sec:failover}

\begin{figure}[htbp]
\centering
\fig{failover-model}
\caption{Failover at the model-node level.}
\end{figure}

\begin{figure}[htbp]
\centering
\fig{failover-inference}
\caption{Failover at the inference-node level.}
\end{figure}

\begin{table}[htbp]
\centering\small
\begin{tabularx}{\linewidth}{@{}YL{6.2cm}@{}}
\toprule
Rule & Value in the code \\
\midrule
Resumes allowed per answer, per level & 2 (\code{MAX_RESUMES}) \\
Frame accepted only if its offset equals the buffer length & Yes; lower offsets dropped, a gap counts as a failure \\
Resume only on servers that say so & \code{Accepted} must carry \code{resume: true} \\
Keep-alive while a model cold-loads & Empty Meta frame every 5\,s \\
A quiet model node is given up on after & 5 to 12\,s, before the app's own 15\,s stall timer \\
\bottomrule
\end{tabularx}
\end{table}

Whoever receives a stream keeps a buffer of exactly what it has passed on: the inference node for
model-node streams, and the app's \code{weirdnet-agent} pool for inference-node streams. On a
reset, a dropped connection or a stall, the receiver reruns routing with the failed node
excluded and sends the same assembled messages, sampling parameters and job id with the next
attempt number, plus its buffer as an assistant prefix. The new engine continues the open
assistant message (vLLM \code{continue_final_message}, llama.cpp assistant prefill) and generates
only what comes after it.

Duplicates are prevented by offsets: every Delta carries its position in the answer, the receiver
appends a frame only if its offset matches the buffer length, and frames from an old attempt are
dropped. Up to two resumes per answer per level are allowed, the final metadata records each
resume, and a client resumes only on a server that says it honours prefixes. Sampling after a
resume is a fresh draw, so the continuation is valid but not bit-identical. If no node can take
over, the app receives the partial answer with a clear error.

Both levels are built. The exit tests kill a model node and then an
inference node mid-answer on separate daemons over loopback with a mock engine, and check that
each reply ends with the exact text.

% ---------------------------------------------------------------------------
\subsection{Context sync}\label{sec:context}

\begin{figure}[htbp]
\centering
\fig{context-sync}
\caption{How a conversation reaches the user's other devices.}
\end{figure}

A request that names a conversation does not wait for the next pass: the inference node first
pulls any missing records from the caller, which always has them, and then assembles the prompt.
Because an answer is stored under its job id, a retry or resume rewrites the same record instead
of adding a second answer.

Conversations are replicated with \code{weirdnet-sync} among inference nodes and app nodes, so any
inference node can continue any conversation and routing never depends on a sticky node. Model
nodes are not in the sync set; they receive fully assembled messages per request and keep nothing
afterwards beyond their local KV cache.

\begin{table}[htbp]
\centering\small
\begin{tabularx}{\linewidth}{@{}lY@{}}
\toprule
Item & Synced? \\
\midrule
Messages & Yes, as text; an assistant message is written once, when it completes \\
System prompt, model, settings, title & Yes, last-writer-wins registers \\
Attachments & Reference synced; bytes fetched on demand (design) \\
KV cache & No; stays on the node that computed it \\
Semantic cache entries & No; pooled separately by \code{weirdnet-semcache} (design) \\
\bottomrule
\end{tabularx}
\end{table}

Records are sealed under the roster's sync key (XChaCha20-Poly1305), which rotates when a device
is unlinked, so a revoked device cannot read new context even through a relay. Deleting a
conversation writes a tombstone that every device honours. Ymir records use their own
\code{ymir.} namespace inside the identity's sync group, which comes from \code{weirdnet-sync}; if
the same identity also runs Morph, another Weirdnet app, the group is shared with it, but each
device opts into only the record types it needs, so a GPU server never holds browsing data.
Bragi's chat history is this same store; there is no separate Bragi history database.

% ---------------------------------------------------------------------------
\subsection{Security and trust model}\label{sec:security}

\begin{figure}[htbp]
\centering
\fig{device-link}
\caption{Linking and unlinking a device.}
\end{figure}

\begin{xltabular}{\linewidth}{@{}L{4.2cm}L{3.6cm}Y@{}}
\toprule
Socket & Bound to & Who can use it \\
\midrule
\endhead
Weirdnet listener & The network & Noise-authenticated peers that pass the roster check \\
OpenAI bridge & 127.0.0.1:8484 only & Local tools with the bridge token and a local Host header \\
MCP over HTTP (optional) & Loopback only & Local MCP clients with the bridge token \\
Engine ports (llama.cpp, vLLM, MLX-LM) & 127.0.0.1 & Ymir, using a random key per launch \\
llama.cpp RPC workers (split groups) & A declared private-link address & The group leader; public and CGNAT addresses refused \\
\bottomrule
\end{xltabular}

\begin{rowlist}
  \item[One authorization check] Every serving node runs the roster check per call: the caller must be a linked device of the owner. No roles, memberships, invites or API keys.
  \item[No network listeners besides Weirdnet] The OpenAI bridge, MCP HTTP mode and engine ports bind to loopback only, and engines get a random per-launch key. The Weirdnet listener accepts only Noise-authenticated peers, then the roster check.
  \item[Revocation] Unlinking a device cuts it off everywhere as peers see the roster update and rotates the sync key. Ymir also refuses a revoked model node on connect, on every request and on every capability update.
  \item[Least data per role] Model nodes see only the prompt in hand; context lives on inference and app nodes; capability records and context go only to linked devices, never the DHT.
  \item[Private-link split workers] llama.cpp RPC workers have no auth of their own, so Ymir binds them only to private addresses (loopback, RFC 1918, link-local, ULA) and refuses public and CGNAT ranges.
  \item[Local secrets] Secrets are never inline in config; the bridge token and identity seed are 0600 files and the state directory is 0700.
\end{rowlist}

The trust boundary is the owner's identity. Today, model and inference nodes receive the owner's
root identity seed when they are linked, which means a copied seed is the owner on Weirdnet.
Milestone M10, a cert-only server link, is designed and approved \planned: service nodes would get
a node certificate and the sync keys they need without the seed, so one node can be revoked
without rotating the identity. The Weirdnet crates for it are published, and Ymir's wiring is
next. The current security review was done by the author and is not an independent audit.

% ---------------------------------------------------------------------------
\subsection{Engines and split models}\label{sec:engines}

\begin{table}[htbp]
\centering\small
\begin{tabularx}{\linewidth}{@{}lYl@{}}
\toprule
Platform & Chat models & Optional \\
\midrule
Linux, Windows & llama.cpp, bundled (CPU and Vulkan, or CUDA 13 build) & vLLM on Linux with NVIDIA \\
macOS, Apple Silicon & MLX-LM with MLX weights & llama.cpp via a flag \\
macOS, Intel & llama.cpp, bundled & \\
\bottomrule
\end{tabularx}
\end{table}

Engines run as child processes under a supervisor rather than in-process, so an engine crash
restarts the engine and triggers failover instead of taking down the daemon. The supervisor does
what llama-swap did per machine: swap groups decide which models may be resident together, idle
models unload after a TTL, crashed engines restart with backoff, a model with requests in flight
is drained before it is stopped, and a small stand-in model can answer, flagged, while a large one
cold-starts. A node can also wrap a stock \code{vllm serve} or \code{llama-server} already running
on loopback. Images and audio route only to nodes whose copy of a model takes them. Video is designed
for after v1: first as frames and audio the client samples, within limits each node reports for its
models (never invented defaults), then native video for models that read it \planned.

A model too big for any one machine can be split across a group of model nodes with llama.cpp RPC
(milestone M7, in progress). It is opt-in per model through pool policy, requires a declared
private link of at least 10\,Gb/s by default, and tears the group down if any member is lost. The
design is explicit that splitting is for models that do not fit, not for speed: in the repo's one
recorded test, an 8B model split between a GPU and a CPU worker ran at 8.3 tokens/s against 63.1
tokens/s unsplit on the GPU. vLLM with Ray groups (\texttt{vllm-ray}: a Ray head and members on the private link, \texttt{vllm serve --distributed-executor-backend ray} on the leader) are built and tested with mock \texttt{ray} and \texttt{vllm}; a run on real GPUs is \planned.

\begin{figure}[htbp]
\centering
\fig{split-throughput}
\caption{Qwen3-8B Q4\_K\_M, measured once in the repo's real-engine split test
(\code{real_split.rs}), with a 3:1 tensor split over loopback; about a quarter of the layers ran
on the CPU and every token crossed the RPC link.}
\end{figure}

\begin{figure}[htbp]
\centering
\fig{split-group}
\caption{A split group for one model.}
\end{figure}

\code{ymirctl group up} checks the policy and links, starts the workers first, then the leader,
and rolls back if a step fails. A watcher checks every 5 seconds and takes the whole group down if
any member leaves, drains or crashes, after which the model returns to normal routing. Starting
groups automatically on demand is not built yet \planned.

% ---------------------------------------------------------------------------
\subsection{Batch fan-out, bridge and MCP}\label{sec:batch}

\begin{figure}[htbp]
\centering
\fig{batch-flow}
\caption{\code{ymirctl batch} fan-out.}
\end{figure}

Because the per-device cap answers Busy, a big batch slows itself to what the pool will take
instead of crowding out interactive chat. Today the submitting machine must stay up while the job
runs; handing a job to another node from a synced job record is not built yet \planned.

\begin{figure}[htbp]
\centering
\fig{bridge-path}
\caption{Local tools reach the whole pool through the loopback bridge.}
\end{figure}

The bridge translates \code{/v1/models}, \code{/v1/chat/completions}, \code{/v1/completions} and
\code{/v1/embeddings} into pool calls, maps Busy to 503 with Retry-After, and marks stand-in and
resumed answers with \code{x-weirdnet-*} headers. \code{ymirctl mcp} offers the same pool to MCP
clients over stdio or loopback HTTP; tools that change the pool refuse to run without
\code{confirm: true}.

\code{ymirctl bench} compares every model on the pool, on every node that serves it. By default it
asks each inference node directly over Weirdnet (\code{ymir-core::direct}), pinning the node so the
same model can be compared across machines; \code{--via bridge} goes through the loopback bridge
instead and records the node that answered. Per model it measures time to first token, decode and
prompt speed and total latency (median and p90 after a warm-up), then scores deterministic quality
checks at temperature 0 (arithmetic, JSON, code run in a sandbox, facts, and a harder suite). It runs
one request at a time, skips a busy or loading model, skips a model that is not loaded unless
\code{--cold}, and gives every request its own seed and nonce so no cache flatters the numbers.
Results are saved as JSON and Markdown, and runs can be compared or exported as CSV.
